\documentclass[12pt]{article}
\usepackage[T1]{fontenc}
\usepackage[utf8]{inputenc}
\usepackage{lmodern}
\usepackage{textcomp}
\usepackage{enumitem}
\usepackage{geometry}
\geometry{margin=1in}
\begin{document}
\begin{center}
Answer Key - Science - K-12 Level Assessment 9 \
\small (Answers are ordered exactly as in the question paper.)
\end{center}

\section*{Multiple Choice}
\begin{enumerate}[label=\arabic*.]
\item \textbf{Answer:} D
\\ \textit{Options:} A) change the experiment.; B) keep their opinions to themselves.; C) find out what other scientists think about the results.; D) repeat the experiment several times and compare results.
\\ \textit{Note:} Repeating the experiment several times ensures that the results are reliable and helps identify any inconsistencies or variables that may have influenced the initial findings.
\item \textbf{Answer:} C
\\ \textit{Options:} A) water; B) salt; C) light; D) oxygen
\\ \textit{Note:} Large blooms of algae on the surface of a lake block sunlight from penetrating through the water, preventing light from reaching the bottom.
\item \textbf{Answer:} C
\\ \textit{Options:} A) balance; B) hot plate; C) stopwatch; D) thermometer
\\ \textit{Note:} A stopwatch is the most appropriate tool for measuring time, allowing for accurate tracking of the duration it takes for a cup of water to reach its boiling point.
\item \textbf{Answer:} D
\\ \textit{Options:} A) their sleek skin.; B) swimming very fast.; C) their body shape.; D) traveling alone.
\\ \textit{Note:} The correct answer is "traveling alone" because dolphins are social animals that typically thrive in groups, known as pods, which provide benefits such as protection, hunting efficiency, and social interaction, making solitary behavior a disadvantage for their survival in the ocean.
\item \textbf{Answer:} B
\\ \textit{Options:} A) current wind conditions; B) type of clouds in the sky; C) percent humidity in the air; D) amount of rainfall
\\ \textit{Note:} The type of clouds in the sky, particularly the presence of cumulonimbus clouds, indicates atmospheric instability and the likelihood of severe weather conditions that can lead to tornado formation.
\end{enumerate}

\section*{True / False}
\begin{enumerate}[label=\arabic*.]
\setcounter{enumi}{5}
\item \textbf{Answer:} False
\\ \textit{Options:} A) True; B) False
\\ \textit{Note:} The statement is false because disturbance encompasses both immediate changes in an ecosystem due to events like glacier melting and the subsequent processes that lead to the establishment of new communities.
\item \textbf{Answer:} False
\\ \textit{Options:} A) True; B) False
\\ \textit{Note:} Keeping the order in which players throw constant can introduce bias and affect performance due to factors like fatigue or pressure, making it an unreliable measure of who can throw the farthest.
\item \textbf{Answer:} False
\\ \textit{Options:} A) True; B) False
\\ \textit{Note:} The statement is false because substances that start off as liquids do not expand and melt when heat is applied; instead, they typically expand and may boil or vaporize, while solids expand and melt when heated.
\item \textbf{Answer:} True
\\ \textit{Options:} A) True; B) False
\\ \textit{Note:} Evaporation occurs at the highest rate in estuaries because their warm temperatures promote faster water molecules escaping into the air, while high salinity levels can increase the concentration of solutes, further enhancing the evaporation process.
\item \textbf{Answer:} True
\\ \textit{Options:} A) True; B) False
\\ \textit{Note:} Electricity generates heat through resistance in heating elements, allowing ovens to cook food effectively.
\end{enumerate}

\section*{Long Form Response}
\begin{enumerate}[label=\arabic*.]
\setcounter{enumi}{10}
\item \textbf{Answer:} Fossil fuels, including coal and natural gas, originate from the remains of ancient plants and animals that lived millions of years ago. Coal forms primarily from the accumulation of plant material in swampy environments, where it is buried under layers of dirt and rock, leading to heat and pressure that transform it into coal over time. Natural gas, on the other hand, is often formed from the remains of tiny marine organisms that also underwent heat and pressure, typically in sedimentary rock formations. Both fossil fuels require specific geological conditions, such as deep burial and the right temperature ranges, to develop effectively. These processes take millions of years, making fossil fuels non-renewable resources that we rely on today.
\\ \textit{Note:} correct because it accurately describes the biological origins of fossil fuels, detailing how coal forms from plant material in swampy areas and how natural gas originates from marine organisms, both processes involving significant heat and pressure over geological time.
\item \textbf{Answer:} When the farmer gathers the crops, the reduction of hiding places for mice can significantly impact the hunting behavior of owls. With fewer places to hide, mice may become more vulnerable to predation, leading to an increase in owl hunting success. This change can affect the local ecosystem by altering the balance between predator and prey populations. If the owl population increases due to the abundance of food, it may help control the mouse population, preventing overpopulation and potential crop damage. Overall, the farmer's actions can create a ripple effect in the ecosystem, influencing various species interactions.
\\ \textit{Note:} correct because it explains how the reduction of hiding places for mice increases their vulnerability to predation by owls, which can lead to changes in predator-prey dynamics and ultimately impact the local ecosystem balance.
\end{enumerate}

\vfill
\noindent\textit{End of Answer Key}
\end{document}